\usepackage{fontspec}
\setmainfont{Georgia}[Ligatures=TeX]
\setsansfont{Georgia}
\setmonofont{Fira Code}[Scale=0.85]
\usepackage{xcolor}
% A4 con márgenes de 2,5 cm: el \textwidth por defecto de LaTeX (345 pt =
% 122 mm) dejaba 44 mm de margen y las imágenes salían espichadas.
\usepackage[a4paper,margin=2.5cm,headheight=14pt,headsep=7mm,footskip=12mm]{geometry}
\usepackage{fancyhdr}
\usepackage{titlesec}
\usepackage{listings}
\usepackage{float}

\definecolor{codebg}{RGB}{245,245,245}
\definecolor{codeframe}{RGB}{110,110,120}
\definecolor{codekeyword}{RGB}{0,90,180}
\definecolor{codecomment}{RGB}{0,128,0}
\definecolor{codestring}{RGB}{180,40,40}

\definecolor{consejocolor}{RGB}{0,110,90}
\definecolor{consejobg}{RGB}{235,248,243}
\definecolor{bugcolor}{RGB}{200,90,0}
\definecolor{bugbg}{RGB}{253,241,230}

\lstdefinestyle{javastyle}{
  basewidth=0.5em,
  backgroundcolor=\color{codebg},
  frame=single,
  rulecolor=\color{codeframe},
  framerule=0.6pt,
  xleftmargin=16pt,
  framexleftmargin=16pt,
  framexrightmargin=0pt,
  fillcolor=\color{codebg},
  basicstyle=\ttfamily\footnotesize\color{black},
  keywordstyle=\color{codekeyword}\bfseries,
  commentstyle=\color{codecomment}\itshape,
  stringstyle=\color{codestring},
  showstringspaces=false,
  breaklines=true,
  numbers=left,
  numberstyle=\ttfamily\tiny\color{gray},
  numbersep=8pt,
  tabsize=4,
  aboveskip=12pt,
  belowskip=12pt,
  linewidth=1\linewidth,
  captionpos=b,
  keepspaces=true,
}

\lstset{style=javastyle}

\setlength{\parindent}{0pt}
\setlength{\parskip}{6pt}

\pagestyle{fancy}
\fancyhf{}
\fancyhead[LE,RO]{\small\itshape Aprende Java Jugando}
\fancyfoot[C]{\thepage}
\renewcommand{\headrulewidth}{0.4pt}

% Títulos en Georgia Bold y en negro: con \sffamily salían en DejaVu Sans
% (mientras el cuerpo es Georgia) y en azul, que no coincide con los Canva.
\titleformat{\section}
  {\Large\bfseries}
  {\thesection}{0.5em}{}
\titlespacing*{\section}{0pt}{0pt}{0.6em}
\newcommand{\sectionbreak}{\clearpage}
\titleformat{\subsection}
  {\large\bfseries}
  {\thesubsection}{0.5em}{}

\usepackage{hyperref}
\hypersetup{
  colorlinks=true,
  linkcolor=blue!60!black,
  urlcolor=blue!60!black,
  pdftitle={Aprende Java Jugando - El Triki},
  pdfauthor={Carlos Alberto Robayo Delgado}
}